\chapter{3D Rendering: Models and Meshes}
\label{ch:models-meshes}

\section{The runtime hierarchy}

\cnaclass{Model} owns a \cnaclass{ModelBoneCollection} and a \cnaclass{ModelMeshCollection};
each \cnaclass{ModelMesh} owns a \cnaclass{ModelMeshPartCollection} and a
\cnaclass{ModelEffectCollection}; each \cnaclass{ModelMeshPart} references a shared
\cnaclass{VertexBuffer}/\cnaclass{IndexBuffer} pair and an \cnaclass{Effect}. This is the
same shape as real XNA, and \cnaclass{Model}'s \emph{runtime} API --- as opposed to how a
model actually gets loaded, \S\ref{sec:model-loading} below --- is fully audited against FNA
with zero known open gaps as of the most recent fix (a \texttt{rootBoneIndex} correction).
\cnaclass{Model} exposes \texttt{CopyAbsoluteBoneTransformsTo},
\texttt{CopyBoneTransformsFrom}/\texttt{To}, and a convenience \texttt{Draw(world, view,
projection)}; one small, intentionally-kept deviation from FNA is that CNA's
\texttt{Copy*BoneTransforms*} methods loop by \texttt{Bones.Count} rather than trusting the
caller's array length the way FNA does --- strictly safer against an oversized destination
array, never a compatibility problem for correctly-sized callers.

\cnaclass{ModelBone} carries \texttt{Name}, \texttt{Index}, a get/set \texttt{Transform},
its \texttt{Parent}, and a \texttt{Children} collection built via a \texttt{NOXNA}
\texttt{AddChild}. \cnaclass{ModelMesh} carries a get/\emph{set} \texttt{BoundingSphere}
(the setter was added specifically to support the real \texttt{.xnb} content pipeline
described below) and a get/\emph{set} \texttt{ParentBone} --- both properties that only
matter once something actually populates them, which is exactly where the two different
model loaders in this codebase diverge sharply.

\section{Two loaders, two very different stories}
\label{sec:model-loading}

This is the single most important thing to understand about \cnaclass{Model} in CNA, and it
is easy to miss if you only read the runtime class headers: \textbf{there are two separate,
independent \cnaclass{Model} loaders}, and \texttt{ContentManager::Load<Model>()} always
tries the newer, real one first.

\begin{description}
  \item[The real \texttt{.xnb} \texttt{ModelReader}] (\texttt{CNA::Internal::Xnb::ModelReader})
        is wire-compatible with real XNA/MonoGame/FNA compiled model assets. It reconstructs
        the full bone hierarchy, assigns each mesh's real \texttt{ParentBone}, computes a real
        \texttt{BoundingSphere}, and correctly resolves shared \cnaclass{VertexBuffer}/%
        \cnaclass{IndexBuffer}/\cnaclass{Effect} resources across meshes that share them ---
        everything below in this section that describes a gap does \emph{not} apply to this
        loader.
  \item[The older \texttt{.model.json} loose-file loader] (\texttt{ModelTypeReader}) is what
        every gap below describes. It always synthesizes exactly one \cnaclass{ModelBone}
        (index~0), even when the source JSON's own \texttt{"bones"} array lists more --- only
        the first entry's name is ever read, and no \texttt{AddChild} hierarchy is ever built.
        Every mesh's \texttt{ParentBone} is left at its default regardless of what the JSON
        says, even though the runtime constructor has supported an explicit parent-bone
        argument since the fix mentioned above --- the loader simply never calls it.
        \texttt{BoundingSphere} is never computed, staying at a degenerate zero radius, and
        neither \texttt{Model.Tag} nor any \texttt{ModelMesh.Tag} is ever set. Resource
        sharing is also absent: every mesh gets a freshly-allocated
        \cnaclass{VertexBuffer}/\cnaclass{IndexBuffer}, and every part using anything other
        than the plain \texttt{"BasicEffect"} string gets its own freshly-constructed
        \cnaclass{Effect} rather than a shared one --- there is no dependency-graph-aware
        resource deduplication of the kind FNA's own \texttt{ReadSharedResource<T>}
        performs. At the time of writing, this loader has \emph{zero} test coverage of its
        own, confirmed by a repository-wide search.
\end{description}

Whether this distinction matters to you depends entirely on which loader your assets go
through --- a real \texttt{.xnb} asset pipeline (Volume~II's migration chapter) sidesteps
every gap in this section; a hand-authored \texttt{.model.json} asset inherits all of them.

\section{Skinning: two independent systems, not one}

CNA ships two separate, unrelated systems for animated meshes, and conflating them is an
easy mistake to make. \cnaclass{AnimationPlayer} (\texttt{NOXNA}, mirroring Microsoft's own
unshipped XNA ``Skinned Model Sample'' reference code rather than anything in the framework
assembly itself) drives a bone-track timeline: \texttt{StartClip(clip)},
\texttt{Update(TimeSpan, relativeToCurrentTime, loop=true)}, and
\texttt{GetBoneTransforms()}/\texttt{GetWorldTransforms()}/\texttt{GetSkinTransforms()} ---
the last of which is exactly what feeds \cnaclass{SkinnedEffect::SetBoneTransforms}
(Chapter~\ref{ch:stock-effects}). Its skinning data (bone count, hierarchy, bind pose,
inverse bind pose, animation clips) is stashed on a loaded \cnaclass{Model}'s own
\texttt{Tag} --- the same convention the original Microsoft sample used, since real XNA's
\cnaclass{Model} has no dedicated skinning-data property of its own to hold it in.

\cnaclass{MorphTargetEXT} (also \texttt{NOXNA}) is a completely independent, glTF-style
morph-target blending system: CPU-side re-blend plus a full vertex-buffer re-upload on every
call to \texttt{SetMorphWeightsEXT}, a deliberate simplicity-over-throughput tradeoff.
Because glTF's ``weights'' animation channel targets a mesh instance node rather than a bone
index, this system has no relationship to \cnaclass{AnimationPlayer}'s bone timeline at all;
it supports linear, step, and real cubic-spline (Hermite) interpolation, and attaches itself
via a loaded \cnaclass{ModelMeshPart}'s own \texttt{Tag} --- the same attachment convention,
applied to a different object in the hierarchy, for a different purpose.

A third, still more separate system exists for the Avatar subsystem specifically:
\texttt{SkinnedModelEXT} (its own \texttt{.skeleton.bin}/\texttt{.clip.bin} binary formats) is
deliberately not built on \cnaclass{Model}/\cnaclass{ModelBone}/\cnaclass{ModelMesh} at all,
has its own full test coverage, and is covered alongside the rest of the Avatar system in
Volume~II.

\section{Collections, uniformly}

\cnaclass{ModelBoneCollection}, \cnaclass{ModelMeshCollection},
\cnaclass{ModelMeshPartCollection}, and \cnaclass{ModelEffectCollection} all follow the same
shape: integer indexing via \texttt{operator[]}, name-based indexing where it makes sense,
\texttt{Count}, \texttt{Contains}, \texttt{TryGetValue}, and standard iterators --- the same
collection idiom already seen on \cnaclass{GameComponentCollection}
(Chapter~\ref{ch:game-loop}) and \cnaclass{CurveKeyCollection}
(Chapter~\ref{ch:math-core-types}), applied consistently a third time here.
